\begin{textsample}{POS Dim 3 – 2016\_c6 – Score -4.00 – 2016\_c6/enade\_2016\_tec\_est\_cos\_fbccf88b8d.txt}  \label{ex:f3_pos_248}
CAPÍTULO 3 – ANÁLISE TÉCNICA DA PROVA Este capítulo tem por objetivo apresentar o desempenho dos estudantes concluintes de Tecnologia em Estética e Cosmética no Enade/2016.
Para isso, foram calculadas as estatísticas básicas da prova em seu todo ( seção 3.1.1 ), bem como as estatísticas relacionadas aos Componentes de Formação Geral ( seção 3.1.2 ) e de Conhecimento Específico da Área ( seção 3.1.3 ).
Dadas as suas características, foram analisadas, em separado, as questões objetivas ( seção 3.2 ) e as questões discursivas ( seção 3.3 ).
Para as questões objetivas, foram disponibilizados os índices de facilidade e de discriminação ponto-bisserial, também em separado, para os Componentes de Formação Geral ( seção 3.2.1 ) e de Conhecimento Específico ( seção 3.2.2 ).
De cada componente, uma das questões foi escolhida para exemplificar a análise gráfica, relacionando as alternativas escolhidas pelos estudantes ( inclusive o gabarito ) com o número de acertos no componente.
O Anexo I apresenta a íntegra da análise gráfica para todas as questões objetivas.
Para cada uma das questões discursivas, os conteúdos dos tipos mais comuns de respostas dos estudantes são apresentados e comparados com o padrão de respostas esperado ( ver Anexo VIII com o padrão de respostas ).
Tomando como base as duas questões discursivas do Componente de Formação Geral, a seção 3.3.1.6 apresenta comentários sobre a correção das respostas com respeito à Língua Portuguesa.
Nas Tabelas 3.1 e 3.3 são apresentados o tamanho da população inscrita, ausente, presente e percentual de faltosos, respectivamente, por Grande Região e por Categoria Administrativa e Organização Acadêmica da IES.
As demais tabelas apresentam as seguintes estatísticas das notas15 : média do desempenho na prova, erro padrão da média, desvio padrão, nota mínima, mediana e nota máxima para cada um de seus componentes.
Tais estatísticas contemplam o total de estudantes concluintes da Área de Tecnologia em Estética e Cosmética inscritos e presentes à prova do Enade/2016, tendo em vista agregações ou por Grandes Regiões e o país como um todo, ou por Categoria Administrativa e Organização Acadêmica da IES.
Em relação aos gráficos de distribuição de notas, o intervalo considerado foi de 10 unidades, aberto à esquerda e fechado à direita, com exceção do primeiro intervalo, [ 0 ; 10 ], fechado em ambos os extremos.
Para os gráficos de distribuição das notas das 15 Uma definição dessas estatísticas pode ser encontrada no Glossário. 3.1 - CONDIÇÃO DE PRESENÇA E ESTATÍSTICAS BÁSICAS DA PROVA Esta seção apresenta a condição de presença, além das estatísticas selecionadas e histogramas da nota geral ( 3.1.1 ) e de cada componente : Formação Geral ( 3.1.2 ) e Conhecimento Específico ( 3.1.3 ).
São também apresentadas estatísticas selecionadas de subpopulações caracterizadas por Grande Região, Categoria Administrativa e Organização Acadêmica. 3.1.1 - Condição de Presença e Estatísticas Básicas Gerais A Tabela 3.1 apresenta a Condição de Presença na prova, por Grande Região, dos estudantes concluintes de Tecnologia em Estética e Cosmética.
A população total de Inscritos foi de 4.686.
Destes, 4.167 estiveram Presentes, sendo 11,1 \% o índice de não comparecimento.
A Região de maior ausência foi a Norte ( 28,0 \%, ) e a de menor ausência foi a Sul ( 5,0 \% ).
Já a Tabela 3.2 apresenta as Estatísticas Básicas da prova, por Grande Região, dos estudantes concluintes de Tecnologia em Estética e Cosmética.
A Média das notas da prova como um todo ( nas seções seguintes serão analisados os Componentes de Formação Geral e de Conhecimento Específico ) foi 43,2, sendo que os alunos da região Norte obtiveram a Média mais baixa ( 38,1 ), e os da região Sul obtiveram a Média mais alta ( 45,0 ).
As demais médias foram : 40,9 na região Centro-Oeste, 42,9 na região Nordeste e 43,2 na região Sudeste.
O Desvio padrão para o Brasil como um todo foi 12,3, sendo o maior Desvio padrão encontrado na região Centro-Oeste ( 12,7 ) e o menor, na região Norte ( 11,5 ), indicando uma dispersão um pouco menor das notas desta última região.
A região que obteve a maior nota Máxima foi a Sul ( 78,4 ), ao passo que a região que atingiu a menor nota Máxima foi a Nordeste ( 73,6 ).
A Mediana do Brasil como um todo foi 42,6, sendo a maior Mediana obtida na região Sul ( 44,7 ), e a menor obtida no Norte ( 36,3 ).
A nota Mínima do Brasil foi zero, obtida na região Sudeste.
Na região Sul a nota Mínima foi 4,8, na Nordeste foi 5,1, na Norte foi 5,9 e na região Centro-Oeste foi 12,2.
Considerando -se as notas segundo Grande Região, observa -se que existe diferença estatisticamente significativa ao nível de 95 \% entre a menor Média, obtida na região Norte ( 38,1 ), e as médias de todas as demais regiões17.
O comportamento das notas dos estudantes de todo o Brasil pode ser observado no Gráfico 3.1 que apresenta um histograma com a distribuição das mesmas.
Essa é uma distribuição unimodal com a moda no intervalo ( 40 ; 50 ], seguido de perto pelo intervalo ( 30 ; 40 ]. 17 Todas as comparações deste capítulo utilizam os intervalos de 95 \%.
Os erros-padrão da média que possibilitam os testes estão disponíveis nas tabelas.
A Tabela 3.3 apresenta informações referentes à Condição de Presença desagregadas por Categoria Administrativa e por Organização Acadêmica.
Da população total de Inscritos, 4.578 são de IES Privadas e 108 de IES Públicas.
Em relação à Organização Acadêmica, a maior participação foi obtida por estudantes de Universidades ( 1.968 ), seguido de Centros Universitários ( 1.642 ) e de Faculdades ( 1.076 ).
A Categoria Administrativa de maior ausência foi a Pública ( 14,8 \% ), e entre as Organizações Acadêmicas foi a dos Centros Universitários ( 12,4 \% ), que ficou acima da taxa de ausência nacional de 11,1 \%.
Já a Tabela 3.4 apresenta informações referentes à Média da nota final desagregadas por Categoria Administrativa e por Organização Acadêmica.
A Média das notas da prova como um todo foi 43,2.
Em relação à Categoria Administrativa, os estudantes das IES Públicas obtiveram Média mais baixa ( 42,3 ) que a nacional, e os das IES Privadas obtiveram a mesma Média ( 43,2 ), que a Média do Brasil.
Observa -se que não existe diferença estatisticamente significativa entre as médias das notas das IES Públicas e Privadas.
A diferença entre as médias das regiões Sul e Norte ( 6,9 ), a maior e a menor Média, é superior à diferença entre IES Públicas e Privadas ( 0,9 ), caracterizando uma menor diversidade administrativa do que regional.
No tocante à Organização Acadêmica, os Centros Universitários e as Universidades obtiveram Média mais alta que a nacional ( 44,3 e 43,4, respectivamente ).
A Média das Faculdades foi 41,1.
O Desvio padrão para as IES Públicas ( 11,5 ) foi inferior ao do Brasil, e o Desvio Padrão das IES Privadas foi igual ao nacional ( 12,3 ).
Dentre as Organizações Acadêmicas o Desvio Padrão variou de 12,0 ( o menor ) nas Faculdades a 12,5 ( o maior ) nos Centros Universitários.
Constata -se que existe diferença estatisticamente significativa ao nível de 95 \% na média das notas dos estudantes das Faculdades ( a menor das médias dos três tipos de Organização Acadêmica ) e as médias dos estudantes dos outros dois tipos de Organização Acadêmica. 3.1.2 - Estatísticas Básicas no Componente de Formação Geral A Tabela 3.5 apresenta as Estatísticas Básicas em relação ao componente da prova que avalia a Formação Geral dos estudantes concluintes.
Os alunos de todo o Brasil obtiveram desempenho médio de 37,9.
Quanto à variabilidade, o Desvio padrão das notas dos estudantes do Brasil, como um todo, foi 15,3.
A maior Média foi obtida na região Sul ( 39,1 ), e a menor, na região Norte ( 35,3 ).
As demais médias foram : 36,4 na região Nordeste, 37,3 na região Centro-Oeste e 37,9 na região Sudeste.
Já o maior Desvio padrão foi obtido na região Nordeste ( 15,9 ), e o menor, na região Norte ( 13,6 ).
Os demais desvios padrões foram : 15,0 na região Centro-Oeste, 15,2 da região Sudeste e 15,6 na região Sul.
A maior nota no Componente de Formação Geral da prova do Enade/2016 foi 90,4, obtida por, pelo menos, um aluno na região Sul.
A menor nota Máxima foi obtida na região Nordeste ( 77,8 ).
A Mediana do Brasil, como um todo, foi 38,0, sendo a menor Mediana encontrada na região Norte ( 35,0 ) e a maior encontrada na região Sul ( 39,8 ).
A nota Mínima nesta parte foi zero em todas as regiões, sem exceção.
Considerando -se as notas segundo Grande Região, observa -se que existe diferença estatisticamente significativa entre a maior Média das notas do Componente de Formação Geral, obtida na região Sul ( 39,1 ), e a menor Média, obtida na região Norte ( 35,3 ).
O Gráfico 3.2 propicia a avaliação do desempenho dos estudantes do Componente de Formação Geral, a partir do histograma da distribuição das notas correspondentes.
A distribuição é unimodal, com moda em ( 30 ; 40 ], enquanto na prova, como um todo ( Gráfico 3.1 ), a moda foi alcançada no intervalo imediatamente posterior, ( 40 ; 50 ].
Nota -se, ainda que, no Gráfico 3.2, as notas apresentam maior dispersão do que no Gráfico 3.1 ( distribuição das notas da prova ), confirmado pela comparação dos desvios padrões : 12,3 para a nota da prova como um todo e 15,3 para o Componente de Formação Geral.
Na Tabela 3.6 são apresentadas as informações referentes ao desempenho dos concluintes do Componente de Formação Geral, em diferentes agregações : Categoria Administrativa e Organização Acadêmica.
Observa -se que não existe diferença estatisticamente significativa entre as médias dos tipos de Categoria Administrativa.
A maior Média foi obtida por estudantes de IES Públicas ( 40,1 ), sem diferença estatisticamente significativa da obtida por estudantes de IES Privadas ( 37,8 ).
Considerando -se o tipo de Organização Acadêmica, nota -se que há diferença estatisticamente significativa entre a média das Faculdades ( a menor das médias ) e as médias dos outros dois tipos de Organização Acadêmica. 3.1.3 - Estatísticas Básicas do Componente de Conhecimento Específico A Tabela 3.7 apresenta as Estatísticas Básicas referentes ao Componente de Conhecimento Específico da área de Tecnologia em Estética e Cosmética.
A Média do desempenho dos alunos do Brasil, como um todo, foi 44,9.
A maior Média foi obtida na região Sul ( 47,0 ), e a menor, na região Norte ( 39,1 ).
As demais médias foram : 42,0 na região Centro-Oeste, 44,9 na região Sudeste e 45,1 na região Nordeste.
Quanto à variabilidade das notas, o Desvio padrão do Brasil, como um todo, foi 13,7, sendo o maior Desvio padrão observado na região Centro-Oeste ( 14,4 ), e o menor, na região Norte ( 12,8 ).
Os demais desvios foram : 13,3 na região Sudeste, 13,6 na região Nordeste e 13,7 na região Sul.
A Mediana das notas dos estudantes de todo o Brasil foi 44,2.
A maior Mediana ocorreu na região Sul ( 46,5 ), e a menor, na região Norte ( 37,1 ).
As demais medianas foram : 40,7 na região Centro-Oeste, 44,0 na região Sudeste e 45,2 na região Nordeste.
A nota Máxima do Brasil, como um todo, foi 84,4, sendo obtida por, pelo menos, um aluno na região Sudeste.
As demais notas máximas foram : 79,8 nas regiões Norte e Nordeste, 82,8 na região Sul e 83,4 na região Centro-Oeste.
A nota Mínima foi zero nas regiões Norte e Sudeste, 2,3 na região Sul, 4,1 na região Nordeste e 4,8 na Centro-Oeste.
Observa -se que existe diferença estatisticamente significativa da menor Média das notas do Componente de Conhecimento Específico, da região Norte ( 39,1 ) em relação às demais regiões.
Assim como os Gráficos 3.1 e 3.2, o Gráfico 3.3, apresentado a seguir, permite uma avaliação do desempenho de concluintes de Tecnologia em Estética e Cosmética em relação ao Componente de Conhecimento Específico com um histograma da distribuição das notas correspondentes.
Esta também é uma distribuição unimodal, e o grupo modal é o ( 40 ; 50 ], o mesmo grupo modal da prova como um todo e acima do grupo modal para a Formação Geral.
A Tabela 3.8 apresenta uma comparação dos resultados em relação à Categoria Administrativa e à Organização Acadêmica, agora levando em conta o desempenho de alunos do Componente de Conhecimento Específico da prova.
No que se refere à Organização Acadêmica, a maior Média foi dos Centros Universitários ( 46,1 ), vindo a seguir a das Universidades ( 45,1 ) e, depois, a das Faculdades ( 42,8 ).
O maior Desvio padrão, e acima do valor para o Brasil ( 13,7 ), como um todo, foi o dos Centros Universitários ( 14,0 ).
Os Centros Universitários também obtiveram as maiores nota Máxima ( 84,4 ) e Mediana ( 45,5 ).
As Universidades obtiveram nota Máxima igual a 81,6 e as Faculdades obtiveram nota Máxima igual a 79,6.
As demais medianas foram : 44,5 nas Universidades e 42,0 nas Faculdades.
A nota Mínima foi zero para as Universidades e para as Faculdades, e foi 4,8 para os Centros Universitários.
Observa -se que existem diferenças estatisticamente significativas ao nível de 95 \% no Componente de Conhecimento Específico entre a nota das Faculdades ( a menor das três notas ) em relação às notas das outras Organizações Acadêmicas.
Quanto à Categoria Administrativa, observa -se que não existe diferença estatisticamente significativa entre as médias das IES Públicas ( 43,0 ) e IES Privadas ( 44,9 ).
Neste caso, diferentemente de Formação Geral, a maior Média foi obtida por alunos de IES Privadas de ensino. 3.2 - ANÁLISE DAS QUESTÕES OBJETIVAS Esta seção apresenta estatísticas selecionadas e histogramas das Questões Objetivas de Formação Geral ( 3.2.1 ).
São também apresentadas e comparadas as médias das subpopulações caracterizadas por Grande Região. 3.2.1 - Componente de Formação Geral A Tabela 3.9 apresenta as Estatísticas Básicas relativas às oito questões objetivas do componente da prova que abrange a Formação Geral dos estudantes.
A Média do Brasil foi 30,8.
A menor Média foi encontrada na região Norte ( 27,2 ), e a maior, na região Sul ( 31,9 ).
As demais médias foram : 29,0 na região Centro-Oeste, 29,9 na região Nordeste e 31,2 na região Sudeste.
O Desvio padrão do Brasil foi 17,6, sendo o maior Desvio padrão encontrado nas regiões Sul e Centro-Oeste ( 17,8 ), e o menor, na região Norte ( 16,1 ).
Os demais desvios foram : 17,4 na região Nordeste e 17,5 na região Sudeste.
As medianas ( 25,0 ), e as notas mínimas ( 0,0 ) foram iguais para todas as regiões.
A nota Máxima do Brasil foi 100,0, obtida na região Sul.
As regiões Norte e Centro-Oeste registram nota Máxima 75,0, e as regiões Nordeste e Sudeste, 87,5.
A Tabela 3.10 apresenta o Índice de Facilidade e o Índice de Discriminação ( ponto-bisserial ) para cada uma das questões objetivas do Componente de Formação Geral.
Quanto ao Índice de Facilidade, foram usadas as seguintes cores para diferenciar o nível de dificuldade da questão :  Azul para as questões classificadas com índice Muito fácil ( \&gt ; =0,86 ), verde para as questões classificadas com índice Fácil ( 0,61 a 0,85 ), amarelo para as questões classificadas com Médio ( 0,41 a 0,60 ), vermelho para as questões classificadas com Difícil ( 0,16 a 0,40 ) e roxo para as questões classificadas com Muito difícil ( \&lt ; =0,15 ).
Já quanto ao Índice de Discriminação, foram usadas as seguintes cores para qualificar a questão :  As questões classificadas com índice Fraco receberam a cor vermelho ( \&lt ; =0,19 ), as classificadas com Médio receberam a cor amarelo ( 0,20 a 0,29 ), as classificadas com Bom receberam a cor verde ( 0,30 a 0,39 ) e as classificadas com Muito bom ( \&gt ; =0,40 ) receberam a cor azul.
As questões objetivas do Componente de Formação Geral, segundo o Índice de Facilidade, foram assim avaliadas : das oito questões, nenhuma teve o Índice de Facilidade classificado como Muito fácil.
Uma questão foi tida como Fácil, por ter índice de acertos situado na faixa entre 0,61 e 0,85 ( de 61,0 \% a 85,0 \% de acertos ).
Uma questão foi considerada de dificuldade Médio, situando -se no intervalo entre 0,41 e 0,60 do Índice de Facilidade, ou seja, houve entre 41,0 \% e 60,0 \% de acertos, e outras quatro questões foram classificadas na categoria Difícil, situando -se no intervalo entre 0,16 e 0,40.
Por fim, duas das questões apresentaram menos de 15 \% de acertos, razão pela qual foram classificadas como Muito difícil.
Como já comentado, para análise das questões objetivas, relativas à Formação Geral, segundo o poder de discriminação, utilizou -se o Índice de Discriminação ( ponto-bisserial ).
Nesta análise, as questões foram assim avaliadas : cinco das oito questões apresentaram índice acima ou igual a 0,40 e, assim, foram classificadas com índice Muito bom para esse grupo de estudantes.
Duas questões tiveram Índice de Discriminação Bom, com valor entre 0,30 e 0,39.
Uma questão teve nível Médio e nenhuma teve nível Fraco de discriminação para esse grupo de estudantes.
O Índice de Facilidade variou de 0,10 a 0,69, e o de Discriminação, de 0,28 a 0,48.
As cinco questões com Índice de Discriminação Muito bom figuraram entre os níveis de dificuldade Fácil, Médio e Difícil : uma classificada na categoria Fácil ( a questão 6 ) do Índice de Facilidade, uma, na categoria Médio ( a questão 3 ) e três na categoria Difícil ( questões 1, 2 e 7 ).
Em particular, a questão 2 a que apresentou o maior poder discriminatório, com índice 0,48, e foi considerada Difícil em termos de facilidade, com uma proporção de 0,30 acertos.
O máximo de acertos foi alcançado pela questão 6 com um Índice de Facilidade de 0,69.
A questão de número 5 apresentou um Índice de Facilidade de 0,10, ou seja, um décimo dos estudantes conseguiu resolvê -la, dentro do universo de participantes.
Seu Índice de Discriminação foi Médio ( 0,28 ).
Já a questão 4 obteve Índice de Discriminação Bom, 0,36, porém seu Índice de Facilidade foi Muito difícil ( 0,14 ).
O Gráfico 3.4, para exemplificar, analisa o comportamento da questão de número 2 de Formação Geral.
Trata -se de uma questão considerada difícil, com relação à facilidade e a que obteve o maior Índice de Discriminação dessa parte da prova.
Neste gráfico, cada uma das cinco curvas representa o percentual de respostas em determinada alternativa da questão, em função do número de acertos dos estudantes nessa parte da prova ( Formação Geral/Múltipla Escolha ), antes de possíveis eliminações pelo critério do ponto-bisserial.
A curva em azul corresponde à alternativa A, a correta para esta questão.
Assim, observa -se que entre os estudantes com menor número de acertos, nessa parte do exame, a situação mais frequente foi a escolha da alternativa E ( em vermelho ), incorreta.
Por exemplo, entre os estudantes que acertaram três questões, 37,3 \% a alternativa A ( correta ), 31,2 \% escolheram a alternativa E, 28,0 \%, a alternativa B ( em verde ), 2,3 \%, a alternativa C ( em preto ) e 1,2 \% escolheu a alternativa D ( em roxo ).
Entre os que acertaram três questões de múltipla escolha de formação geral, 0,1 \% deixou a questão em branco e 0,0 \% marcou mais de uma alternativa, o que invalida a questão.
À medida em que o número de acertos aumenta, indicando desempenho melhor nessa parte da prova, aumenta concomitantemente a proporção de estudantes que selecionaram a alternativa correta A, atingindo 100 \% para os estudantes com 7 ou 8 acertos.
Essa análise permite verificar como a questão discriminou os grupos de desempenho, justificando o alto índice obtido na questão.
Cumpre notar que não é possível inferir deste gráfico nem o índice de facilidade ( que seria uma média da proporção ponderada pela quantidade de alunos com cada uma das notas ), nem o índice de discriminação ponto-bisserial, por razão equivalente.
No caso extremo no qual a grande concentração dos acertos dos alunos fosse abaixo de 4, o índice de facilidade seria obrigatoriamente abaixo de 40 \% ( neste exemplo ).
Caso a concentração fosse em 6 acertos ou mais, o índice seria obrigatoriamente acima de 60 \%.
Os gráficos relativos às demais questões de Formação Geral constam do Anexo I. 3.2.2 - Componente de Conhecimento Específico A Tabela 3.11 apresenta as Estatísticas Básicas em relação às questões objetivas do Componente de Conhecimento Específico da prova de Tecnologia em Estética e Cosmética por Grande Região.
A Média do Brasil deste componente foi de 48,0.
Nessa parte da prova uma questão foi anulada pela Comissão Assessora da Área, a questão 32.
Além disso, como será discutido ainda nessa seção, cinco questões objetivas do Componente de Conhecimentos Específicos não foram usadas no cômputo das notas por terem sido descartadas pelo critério do ponto-bisserial.
Assim, as notas dizem respeito ao resultado de 21 das 27 questões objetivas dessa parte da prova.
A menor Média foi observada na região Norte ( 42,2 ), e a maior, na região Sul ( 50,1 ).
O Desvio padrão do Brasil foi 14,6, sendo o menor Desvio padrão encontrado na região Norte ( 13,8 ), e o maior, na região Centro-Oeste ( 15,5 ).
A Mediana do Brasil foi 47,6, a mesma encontrada nas regiões Nordeste, Sudeste e Sul.
Nas regiões Norte e Centro-Oeste, a Mediana foi 42,9.
A nota Máxima da prova foi 90,5, obtida nas questões objetivas do Componente de Conhecimento Específico, por, pelo menos, um aluno na região Centro-Oeste.
Na região Norte, a nota Máxima da prova foi 81,0, e nas demais regiões, foi 85,7.
Nas regiões Norte, Sudeste e Sul a nota Mínima foi zero, e foi 4,8 nas outras duas regiões.
A Tabela 3.12 apresenta os Índices de Facilidade e Discriminação ( ponto-bisserial ) das questões objetivas do Componente de Conhecimento Específico para os estudantes de Tecnologia em Estética e Cosmética.
Para facilitar a diferenciação das questões, usaremos as mesmas cores da Tabela 3.10 para as diferentes classificações dos Índices de Facilidade e de Discriminação.
A partir do Índice de Facilidade obtido, pode -se concluir que mais de um terço das questões objetivas da prova foi considerado Difícil ou Muito difícil : das 26 questões válidas ( a questão 32 foi anulada pela CAA ), 10 foram classificadas como Difícil ou como Muito difícil.
Duas questões foram classificadas como Muito fácil, oito questões foram classificadas como Fácil, e outras seis consideradas como Médio.
Já quanto aos índices de discriminação das questões objetivas do Componente de Conhecimento Específico da prova, tem -se como resultado a seguinte classificação : 13 das 26 questões foram consideradas como boas, enquanto nenhuma delas tive Índice de Discriminação Muito bom.
Assim, para metade das questões – 13 em 26 – o Índice de Discriminação foi Bom.
Dentre as demais, oito delas foram classificadas como Médio e outras cinco como Fraco, sendo 13, por conseguinte, a quantidade de questões nos dois patamares mais baixos de discriminação.
Constata -se, assim, que a prova – no que se refere ao Componente de Conhecimento Específico – possuía razoável capacidade de discriminar entre aqueles que dominam ou não o conteúdo.
A questão de número 33 foi a mais difícil dentre as 26 questões específicas válidas, com baixo Índice de Facilidade 0,14, apenas 14 \% de acertos.
Essa questão apresentou poder discriminatório igualmente baixo, 0,19, o que comprova ter sido esta questão a mais difícil para os estudantes.
Destaca -se, também, a questão 27, com Índice de Facilidade 0,70, o que, em termos percentuais, corresponde a 70 \% de estudantes que responderam acertadamente.
Entretanto, o seu de Índice de Discriminação foi baixo, 0,13.
Tais questões foram, portanto, pelo critério ponto-bisserial, consideradas inadequadas.
Por isso, as questões 27 e 33, além das questões 10, 12 e 24, foram eliminadas do cômputo da nota final.
A título de exemplo das análises do comportamento das questões objetivas, o Gráfico 3.5 analisa a questão 20 do Componente de Conhecimento Específico.
Esta questão foi considerada, pelas respostas dos estudantes avaliados na prova, como uma questão com nível de facilidade Médio, 0,49, ou seja, 49 \% dos estudantes assinalaram acertadamente a opção D, correspondente ao gabarito.
Seu Índice de Discriminação foi igual a 0,39, classificado como Bom, também sendo esta questão a que apresentou o maior valor.
Neste gráfico, cada uma das cinco curvas representa o percentual de respostas em determinada alternativa da questão 20, em função do número de acertos dos estudantes nessa parte da prova, antes de possíveis eliminações de questões pelo critério do ponto-bisserial.
A alternativa correta D, representada no gráfico pela curva em roxo, foi escolhida em maiores proporções pelos alunos com desempenho melhor nessa parte da prova.
Já as alternativas incorretas, também denominadas distratores, foram selecionadas, principalmente, por aqueles com notas mais baixas.
Neste caso também, a soma não é sempre 100 \% por causa das questões não respondidas ou com mais de uma opção marcada.
Aqueles com nota zero, na sua quase totalidade deixaram esta questão em branco ou marcaram mais de uma alternativa, comportamento considerado inválido.
A proporção de alunos que selecionou a resposta correta D aumenta gradativamente, chegando a atingir 100 \% para 23 acertos, enquanto a proporção dos que escolheram alternativas incorretas decai, a partir de oito acertos, como função do número de acertos nessa parte da prova.
Os gráficos relativos às demais questões do Conhecimento Específico constam do Anexo I. 3.3 - ANÁLISE DAS QUESTÕES DISCURSIVAS Esta seção apresenta estatísticas selecionadas e histogramas das Questões Discursivas de Formação Geral ( 3.3.1 ) e Conhecimento Específico ( 3.3.2 ).
São também apresentadas e comparadas as médias de subpopulações caracterizadas por Grande Região. 3.3.1 - Componente de Formação Geral As análises dos resultados de desempenho dos estudantes de Tecnologia em Estética e Cosmética nas duas questões discursivas relativas à Formação Geral encontram -se na Tabela 3.13 e no Gráfico 3.6.
Na Tabela 3.13, observa -se que a nota Média nesse conjunto de questões foi acima da obtida nas objetivas.
Os estudantes, de todo o Brasil, obtiveram, em Formação Geral, Média 30,8 nas questões objetivas e 48,5 nas questões discursivas.
Pode -se notar também que o Desvio padrão nesse conjunto de questões foi superior ao obtido nas objetivas, 17,6 nas questões objetivas e 22,6 nas questões discursivas.
A maior Média foi obtida na região Sul ( 49,9 ), e a menor, na região Nordeste ( 46,1 ).
A Mediana de todo o Brasil, neste componente, foi 54,0, a mesma obtida nas regiões Sudeste e Centro-Oeste.
Na região Sudeste, a Mediana foi maior, 56,0, e nas regiões Note e Nordeste, foi menor que o valor para o Brasil, 53,0 e 52,0, respectivamente.
A nota Máxima ( 97,0 ) foi obtida na região Norte, sendo 88,5 a nota Máxima na região Sudeste, 91,5 na região Centro-Oeste, 92,0 na região Sul e 93,0 na Nordeste.
A nota Mínima ( 0,0 ) foi a mesma em todas as regiões do Brasil.
O Gráfico 3.6 representa a distribuição das notas nas questões discursivas do Componente de Formação Geral.
A moda desta distribuição ocorre no intervalo ( 60 ; 70 ].
Destaca -se também o intervalo [ 0 ; 10 ] com distribuição bem próxima a 10 \% do total de notas, sendo que no intervalo [ 0 ; 10 ] inclui -se, além da nota zero, a frequência de alunos que deixaram este tipo de questão em branco.
Na sequência, os resultados verificados para cada uma das questões discursivas de Formação Geral serão apresentados, estabelecendo -se relações com os conteúdos abordados em cada uma delas.
Os comentários da Banca de docentes corretores, a respeito do observado na correção das respostas dos estudantes, suas impressões e conclusões serão apresentados junto à análise de cada questão.
Cumpre esclarecer que, tendo em vista que as questões discursivas de Formação Geral são padronizadas, ou seja, constam de todas as provas, os comentários da Banca são os mesmos para todas as carreiras acadêmicas, sendo direcionados a todos os estudantes que participaram do Enade/2016.
A seguir, serão analisados os desempenhos da Área de Tecnologia em Estética e Cosmética nas duas questões discursivas de Formação Geral do Enade/2016, comparando -se os resultados obtidos com comentários para cada questão. 3.3.1.1 - Análise de Conteúdo da Questão Discursiva 1 do Componente de Formação Geral Os dados de Tecnologia em Estética e Cosmética, obtidos a partir das respostas à questão 1, encontram -se na Tabela 3.14 e no Gráfico 3.7.
Nessa questão – de melhor desempenho dentre as duas de Formação Geral – os alunos, de todo o Brasil, tiveram Média 50,6.
A maior Média para a questão 1 foi obtida na região Sul ( 52,3 ), e a menor, na região Nordeste ( 49,0 ).
Quanto à variabilidade das notas, o Desvio padrão de todo o Brasil foi 32,7.
O menor Desvio padrão foi obtido na região Centro-Oeste ( 32,0 ), e o maior Desvio padrão foi obtido na região Nordeste ( 34,9 ).
As medianas de todas as regiões e a do Brasil como um todo foram iguais ( 60,0 ).
As notas máximas e mínimas da questão discursiva 1 também foram as mesmas para todas as regiões do Brasil, respectivamente, 100,0 e 0,0.
O Gráfico 3.7 mostra a distribuição das notas na questão discursiva 1 do Componente de Formação Geral.
Observa -se que a maior frequência corresponde aos alunos que obtiveram nota dentro do intervalo ( 70 ; 80 ], correspondendo à moda da distribuição.
Destacam -se, também, o percentual de alunos que deixaram a questão em branco, aproximadamente 17,5 \% do total. 3.3.1.2 - Comentários sobre a correção de Conteúdo das respostas à Questão Discursiva 1 A questão discursiva 1 alinha -se aos conteúdos programáticos previstos nas diretrizes curriculares, quanto às habilidades e às competências para a atualização permanente e aos conhecimentos sobre a realidade brasileira, mundial, tangenciando outras áreas do conhecimento.
O enunciado é claro e indica um tema de extrema relevância, qual seja, as migrações humanas e todo o complexo de problemas que esse movimento causa, quer na origem, quer no seu ponto de destino.
A questão é extremamente atual, maximizada pelos acompanhamentos midiáticos que espetacularizam cenas dramáticas do processo de fuga de massas humanas das guerras, como a morte de imigrantes.
Além do aspecto internacional, deve se ter em conta que o Brasil sofreu vários influxos migratórios, desde o forçado, pela transferência de escravos da África, até o incentivado, dos italianos, alemães e outros europeus para trabalhar na agricultura.
Isso sem considerar a própria colonização portuguesa originária e as migrações posteriores dos mesmos portugueses.
Posteriormente, o Brasil tornou -se um país de emigrantes que foram aportar a outros locais na busca de melhor situação econômica.
Muitos migraram para a América do Norte e para a Europa, principalmente.
Mais recentemente, por força do crescimento econômico pátrio, novas levas de imigrantes desembarcaram no Brasil, muitos deles provenientes do Haiti, embora tenhamos recebidos vizinhos latinos, como bolivianos, paraguaios, chilenos, argentinos, venezuelanos, peruanos, colombianos que também vieram em busca de uma vida melhor.
Como se verifica, a temática indicada no enunciado pertence à realidade conhecida amplamente pelos estudantes, pois, também divulgada por todos os meios de comunicação.
Trata -se de tema geral que, normalmente, não se inclui nas grades curriculares superiores, mas é de conhecimento geral e relacionado aos direitos humanos, este sim, ponto de integração às diretrizes curriculares.
A linguagem utilizada foi clara, indicando a existência dos fatos vinculados aos movimentos migratórios.
Por outro lado, o texto demanda uma reflexão sobre a realidade dos indivíduos que são compelidos ao afastamento dos seus lares por motivos externos, como guerras, conflitos civis, religiosos, perseguições das mais variadas, questões econômicas.
Como se trata de texto contendo uma ampla gama de possibilidades, é natural que ocorram interpretações diversas das esperadas.
Do mesmo modo, a contemporaneidade do tema encaminha o enunciado para ser qualificado como questão de média dificuldade, em comparação com as demais questões apresentadas.
O padrão de resposta indicou possíveis interpretações do texto.
No que concerne ao grau de profundidade, mostrou a necessidade de conhecimento de vários itens pertinentes à migração e que são refletidos nas várias possibilidades de resposta previstas pelo padrão.
Pela análise das provas realizadas, constatou -se que muitos foram os que identificaram os itens indicados no padrão de resposta.
Essa circunstância poderia ser justificada, tendo em vista que o tema vem sendo bastante difundido na mídia social, gerando elementos naturais para a resposta, bem como da história do país, que, intimamente, vincula -se ao processo migratório, quer ativo, remetendo pessoas para outros países, quer passivo, na recepção de pessoas.
Essa característica também permite responder aos termos postos quanto aos Estados que recebem os migrantes.
Assim, a indicação dos aspectos positivos e negativos também é bastante difundida.
As respostas medianas e boas apreenderam o sentido do enunciado, e indicaram os motivos expressos no padrão de respostas, como os ensejadores das imigrações, notadamente, para a Europa, que seria o foco do texto.
Por outro lado, também identificaram os problemas dos imigrantes, no mesmo padrão indicado para as respostas desejadas.
Alguns concluintes procuraram exemplos de comparação na realidade brasileira, como os clássicos imigrantes italianos, alemães e de outras nacionalidades, bem como os atuais haitianos e demais povos da América Latina.
Uma minoria não apreendeu o sentido do enunciado e, somente, apresentou elementos da realidade nacional.
Esses foram avaliados no conjunto de provas fracas.
As respostas demonstraram que o concluinte está atualizado, em regra, com o tema da imigração, indicando os fundamentos dessa movimentação populacional que, nos últimos tempos, tem se agravado pela incidência de conflitos armados em vários países, ensejadores de reflexos na segurança dos habitantes desses locais, bem como o surgimento de focos de doenças e ausência dos serviços públicos básicos, além da fome, decorrente da impossibilidade de produzir alimentos ou mesmo comerciá -los.
Do ponto de vista dos Estados que recebem o fluxo de imigrantes, também restou clara a percepção das dificuldades vinculadas à ausência de estrutura quanto aos serviços públicos ( saúde, educação, segurança, moradia, alimentação ) para receber um número expressivo de imigrantes.
Em época de recessão econômica, a presença de estrangeiros acarreta a competição pela busca de empregos, gera movimentos xenofóbicos, bem como choques culturais e políticos.
Deve -se ter em mira que esses movimentos têm um viés de insegurança, pelos países receptores, diante da crise causada por movimentos considerados terroristas, embora a imensa maioria não tenha qualquer vínculo com esses movimentos.
Os recentes atentados na Europa indicam que os seus autores têm origem em comunidades de imigrantes, embora muitos já tenham nascido nos países da comunidade europeia.
Essas características negativas, no entanto, paradoxalmente, convivem com aspectos positivos, o que encaminha a questão para a necessidade dos jovens imigrantes, que ocuparão postos de trabalho não desejados pelos habitantes locais, bem como, compensarão os baixos índices de crescimento populacional.
Existe outro reflexo positivo que está relacionado ao sistema previdenciário, que terá novos contribuintes nos trabalhadores estrangeiros que se integram ao sistema.
Esses pontos foram abordados e indicam que existe uma consciência dos concluintes dos cursos universitários quanto aos principais itens relacionados ao tema que foi posto para análise. 3.3.1.3 - Análise de Conteúdo da Questão Discursiva 2 do Componente de Formação Geral A Tabela 3.15 mostra que o desempenho médio dos estudantes na questão discursiva 2 ( média 43,1 ) foi inferior ao obtido na questão discursiva 1 ( média 50,6 ).
A região Centro-Oeste foi aquela cuja Média, nessa questão, foi maior ( 44,8 ), e a de menor Média foi a região Nordeste ( 40,3 ).
Quanto à variabilidade das notas, o Desvio padrão de todo o Brasil foi 23,7, inferior ao obtido na questão discursiva 1 ( 32,7 ).
O maior desvio nessa questão foi obtido na região Nordeste ( 24,6 ), enquanto o menor foi obtido na região Sudeste ( 23,1 ).
A Mediana do Brasil e de quatro das cinco regiões foi 50,0 - a única região com Mediana diferente foi a Nordeste, na qual a Mediana foi 40,0.
As notas máximas ( 100,0 ) foram as mesmas em todas as regiões do Brasil, assim como as notas mínimas ( 0,0 ).
O Gráfico 3.8 mostra a distribuição das notas na questão discursiva 2 do Componente de Formação Geral.
Observa -se que a maior frequência corresponde aos alunos que obtiveram nota dentro do intervalo ( 50 ; 60 ].
Destacam -se, também, os alunos que deixaram a questão em branco e os alunos que obtiveram nota dentro do intervalor ( 30 ; 40 ], caracterizando -se como máximos locais.
Nota -se, ainda, que as notas ficaram menos distribuídas ao longo dos intervalos em comparação à questão discursiva de número 1.
Isto pode ser constatado também pela comparação do desvio padrão das notas da questão discursiva 2 ( 23,7 ) e da questão discursiva 1 ( 32,7 ). 3.3.1.4 - Comentários sobre a correção de Conteúdo das respostas à Questão Discursiva 2 A questão discursiva 2 tem seu conteúdo integrado às Diretrizes Curriculares, uma vez que a proteção aos direitos humanos e às relações de gênero constituem matrizes fundamentais para conhecimento dos profissionais de nível superior.
O enunciado é claro e indica um tema de extrema relevância, qual seja, a análise da questão da mulher submetida a injustas agressões e os seus reflexos pessoais, sociais e familiares.
Também enfrenta a busca pela maior proteção no concernente às agressões, com a instituição de um modelo legal mais repressivo e que foi gerado por uma situação especifica de ataque à mulher, com consequências dramáticas para a vítima, Maria da Penha, cuja lei levou seu nome.
Deve ser recordado que os conflitos vinculados às relações familiares sempre foram circunscritos ao seu âmbito e, quando dele extravasavam, por força de agressões contundentes à integridade física de um dos partícipes da relação, em geral a mulher, raramente redundavam em punições relevantes, sendo comum o pagamento em cestas básicas de alimentos, como sanção.
Após a edição da lei Maria da Penha ( Lei no. 11.340, sancionada em 7 de agosto de 2006 ), além das sanções sofrerem modificações, houve a criação de modelos de proteção e atendimento especial para os casos de violência doméstica, procurando proteger as mulheres agredidas antes de danos maiores, muitas vezes irreversíveis, incluindo aí o feminicídio, também incluído no Código Penal, como circunstância qualificadora do crime de homicídio, e incluí -lo no rol dos crimes hediondos.
A lei Maria da Penha busca realizar, como indicado no seu pórtico, a Convenção sobre a Eliminação de Todas as Formas de Discriminação contra as Mulheres e a Convenção Interamericana para Prevenir, Punir e Erradicar a Violência contra a Mulher.
No seu texto, várias são as formas de violência doméstica e familiar contra a mulher, identificadas como sendo, entre outras : “ a violência física, entendida como qualquer conduta que ofenda sua integridade ou saúde corporal ; a violência psicológica, entendida como qualquer conduta que lhe cause dano emocional e diminuição da autoestima ou que lhe prejudique e perturbe o pleno desenvolvimento ou que vise a degradar ou controlar suas ações, comportamentos, crenças e decisões, mediante ameaça, constrangimento, humilhação, manipulação, isolamento, vigilância constante, perseguição contumaz, insulto, chantagem, ridicularização, exploração e limitação do direito de ir e vir ou qualquer outro meio que lhe cause prejuízo à saúde psicológica e à autodeterminação ; a violência sexual, entendida como qualquer conduta que a constranja a presenciar, a manter ou a participar de relação sexual não desejada, mediante intimidação, ameaça, coação ou uso da força ; que a induza a comercializar ou a utilizar, de qualquer modo, a sua sexualidade, que a impeça de usar qualquer método contraceptivo ou que a force ao matrimônio, à gravidez, ao aborto ou à prostituição, mediante coação, chantagem, suborno ou manipulação ; ou que limite ou anule o exercício de seus direitos sexuais e reprodutivos ; a violência patrimonial, entendida como qualquer conduta que configure retenção, subtração, destruição parcial ou total de seus objetos, instrumentos de trabalho, documentos pessoais, bens, valores e direitos ou recursos econômicos, incluindo os destinados a satisfazer suas necessidades ; a violência moral, entendida como qualquer conduta que configure calúnia, difamação ou injúria. ” ( Conforme texto extraído do artigo sétimo da referida lei ).
Também se buscam medidas de integração entre os vários órgãos públicos, abrangendo repressão, proteção, conscientização, bem como utilização de equipes multidisciplinares, tendo em vista os graves danos, de variados matizes, que redundam das práticas agressivas.
Nessa linha, também existe na lei a previsão de locais de atendimento específico ( centros de atendimento integral e multidisciplinar para mulheres e respectivos dependentes em situação de violência doméstica e familiar ; casas-abrigo para mulheres e respectivos dependentes menores em situação de violência doméstica e familiar ; delegacias, núcleos de defensoria pública, serviços de saúde e centros de perícia médico-legal especializados no atendimento à mulher em situação de violência doméstica e familiar ; programas e campanhas de enfrentamento da violência doméstica e familiar ; centros de educação e de reabilitação para os agressores ).
O Poder Judiciário, também, deve especializar o seu atendimento para aplicar as medidas de restrição ao agressor de forma mais eficiente.
A linguagem utilizada na questão foi clara, indicando a grave violação dos direitos das mulheres injustamente agredidas e a luta pela busca de sanções adequadas aos agressores e de proteção pessoal e para os de sua família.
Ademais, o tema não está circunscrito ao Brasil, mas o país possui um número muito expressivo de atos que compõem esse cenário.
Por outro lado, o texto demanda uma reflexão sobre a realidade, e abarca diversos segmentos da sociedade, pois as situações de agressão não se restringem às camadas mais populares.
Como se trata de texto contendo uma ampla gama de possibilidades, é natural que ocorram interpretações diversas das esperadas.
O padrão de resposta identifica a projeção do elaborador na resposta a ser apresentada pelo estudante.
Está em linguagem correta e indica possíveis interpretações do texto.
Suas opções foram bastante amplas, abrangendo a maior parte das respostas encontradas.
No que concerne ao grau de profundidade, o padrão de resposta indicou a necessidade de conhecimento de lei especial que regula a violência doméstica, o que exige ou qualificação da área de Direito ou que os currículos dos cursos dos participantes desse ciclo do Enade tenham oferecido alguma disciplina introdutória da área jurídica.
Ao analisar as provas realizadas, constatamos que, na primeira parte, houve maior facilidade em responder aos itens indicados no padrão de resposta, pois atinentes, não só à formação profissional, mas, também, ao conhecimento da realidade, pois, infelizmente, a violência doméstica está disseminada na sociedade e nos vários pontos do país.
Mas, verificou -se que as respostas vieram, no mais das vezes, desorganizadas, dificultando a correção.
Apesar disso, restou claro que a maioria esmagadora das respostas confluiu para as indicações do padrão apresentado, sendo relevante indicar que o tema remeteu para a luta em relação à igualdade de gênero, temática candente em especial para as concluintes do sexo feminino.
Diversos textos se apresentaram de forma concatenada, fundamentada e, muitas vezes, lamentando que, nos dias atuais, ainda se discuta a necessidade de proteger as mulheres dessa espécie de violência.
Em muitos casos, também, os concluintes fizeram um liame entre o acesso à educação e a melhora das relações entre homens e mulheres, ao permitir que a conclusão de maior grau de educação formal conduzisse à diminuição do índice de violência contra a mulher.
Foram identificadas respostas consoantes com o padrão apresentado, indicando os itens que nele constam em relação aos reflexos pessoais, na família e na sociedade.
As respostas dos estudantes raramente faziam referência à organização protetiva instituída pela Lei Maria da Penha, como esperado pelo padrão de resposta.
A respeito de tal aspecto o padrão indicava como uma das possibilidades de respostas ( item ‘ b ’ ) : a criação de casas de refúgio ou casa-abrigo, para acolher mulheres vítimas de violência doméstica ; reforço às Delegacias de Atendimento à Mulher ; inclusão da vítima em programas assistenciais do governo, programas de proteção à vítima e à testemunha, transferência de local de trabalho, se a vítima for servidora pública.
Como paradigma, indicaram -se as Delegacias da Mulher que são organizadas, com a quase integralidade dos seus membros de policiais do sexo feminino, e especializadas no atendimento a mulheres submetidas a situações de violência.
Essa preocupação gerencial tem por escopo bloquear os preconceitos que acompanham essas dramáticas histórias de violência, o que gera uma incompreensão e uma ineficiência na rede de proteção.
As respostas demonstram que os concluintes têm ciência do problema, da sua gravidade, do vínculo com a questão de gênero e com a submissão da mulher como efeito da sociedade patriarcal, ainda que tratado de forma superficial, com base na mídia e, algumas vezes, com apoio de experiências pessoais.
Identificou -se uma generalizada compreensão de que o problema está sendo enfrentado e que ainda existe um longo caminho no percurso para a sua resolução. 3.3.1.5 - Análise de Língua Portuguesa das Questões Discursivas do Componente de Formação Geral Os dados de Tecnologia em Estética e Cosmética, obtidos a partir das respostas às questões discursivas do Componente de Formação Geral, no que tange à Língua Portuguesa, encontram -se na Tabela 3.16 e no Gráfico 3.9.
Nesse aspecto, os alunos, de todo o Brasil, tiveram Média 55,0.
A maior Média com respeito à Língua Portuguesa foi obtida na região Sul ( 56,4 ), e a menor, na região Nordeste ( 52,0 ).
Quanto à variabilidade das notas, o Desvio padrão de todo o Brasil foi 22,1.
O menor Desvio padrão foi obtido na região Centro-Oeste ( 21,1 ) e o maior Desvio padrão foi obtido na região Nordeste ( 23,3 ).
A Mediana das notas de Língua Portuguesa foi 62,5 para o Brasil como um todo e também para a região Sudeste.
A maior Mediana com respeito à Língua Portuguesa foi obtida na região Sul ( 65,0 ), e a menor, nas regiões Norte, Nordeste e Centro-Oeste ( 60,0 ).
A nota Máxima para todo o Brasil foi de 90,0, com, pelo menos, um aluno tirando essa nota nas regiões Sudeste e Centro-Oeste.
Além disso, a nota Mínima foi zero em todas as regiões do país.
O Gráfico 3.9 mostra a distribuição das notas de Língua Portuguesa do Componente de Formação Geral.
Observa -se que a maior frequência ( em torno de 33 \% ) corresponde aos alunos que obtiveram nota no intervalo ( 60 ; 70 ].
Destacam -se, também, os alunos que deixaram ambas as questões em branco, representando quase 10 \% do total, caracterizando -se como máximo local. 3.3.1.6 - Comentários sobre a correção das respostas de Formação Geral com respeito à Língua Portuguesa Na avaliação do desempenho linguístico dos participantes foram consideradas quatro competências distintas, de modo a permitir um mapeamento detalhado do domínio dos recursos disponíveis na Língua Portuguesa para a comunicação escrita formal.
Nessa situação comunicativa, as exigências para o desenvolvimento desse tipo de texto dizem respeito : à adequação da seleção vocabular, ao desenvolvimento do conteúdo, à estruturação sintática dos períodos, à organização lógica das ideias, à utilização de procedimentos de encadeamento textual e de referenciação, à obediência às exigências morfossintáticas próprias da modalidade escrita da norma-padrão, ao respeito às regras ortográficas e às regras de acentuação gráfica.
A avaliação foi realizada considerando -se os textos das respostas às duas questões de Formação Geral e os resultados obtidos em cada uma delas, como seria de se esperar, foram semelhantes.
As competências avaliadas podem ser sintetizadas nos aspectos descritos a seguir.
Aspectos textuais - estruturação textual condizente com o gênero discursivo e o modo de organização textual expositivo/argumentativo adequado ao gênero – essa competência envolve : a estruturação sintática condizente com o padrão da modalidade escrita formal da Língua Portuguesa de modo a garantir a clareza necessária ; a distribuição do conteúdo do texto em parágrafos, de modo a garantir a sua organização temática ; a utilização de operadores discursivos que contribuam para a progressão temática do texto, estabelecendo relações lógicas entre as ideias apresentadas, tanto do ponto de vista intrafrasal, como do interfrasal ; a utilização de procedimentos de referenciação lexical e pronominal que permitam a retomada de referentes textuais ; o respeito às regras de pontuação como fator de estruturação do período.
Aspectos morfossintáticos - domínio dos diferentes aspectos morfossintáticos próprios da modalidade escrita formal da norma-padrão da Língua Portuguesa – essa competência envolve : a concordância nominal, a concordância verbal, a regência nominal, a regência verbal, a flexão nominal, a flexão verbal, a correlação entre os tempos verbais, a colocação pronominal e a utilização de sinais de pontuação que contribuam para a organização lógica da frase e do texto.
Aspectos vocabulares - seleção vocabular adequada à modalidade escrita formal da Língua Portuguesa, exigida pela situação comunicativa – essa competência envolve a precisão na seleção/utilização do vocabulário relacionado à temática solicitada pela questão ; a ausência de marcas da oralidade, como termos de sentido muito genérico ( “ coisa ”, “ negócio ”, “ você ” ) e termos de registros mais informais ( como gírias, jargões, frases feitas, ditados populares, termos regionais ).
Assim, espera -se que o participante respeite a adequação vocabular não usando gírias ou expressões coloquiais, evite repetição desnecessária de palavras e utilize um vocabulário mais formal, como solicitado por um texto dissertativo.
Aspectos ortográficos - respeito às convenções ortográficas da norma-padrão da Língua Portuguesa – essa competência envolve o domínio das regras de acentuação gráfica e da grafia padrão das palavras ( com ausência de abreviaturas próprias da linguagem da internet ), de acordo com as convenções estabelecidas pela legislação em vigor e consubstanciadas no Vocabulário Ortográfico da Língua Portuguesa, editado pela Academia Brasileira de Letras. ( Vale notar que, nessa edição de 2016, não será aceita a legislação anterior, no caso das regras relativas ao uso do hífen e da acentuação gráfica ).
A escolha dessas competências para subsidiar o processo de avaliação apoia -se na concepção de que, no desempenho dos graduandos, a modalidade escrita tem apresentado uma intensa simplificação, originada no padrão da modalidade oral da Língua Portuguesa.
No caso do texto de base dissertativa, inscrito em um registro formal, a distância entre as duas modalidades é ainda maior, o que provoca situações de hipercorreção ( desvios provocados pela incorporação indevida de uma regra da norma-padrão ) e de truncamentos sintáticos ( estruturas frasais incompreensíveis devido à complexidade sintática própria da modalidade escrita ).
Para efeito de pontuação, o critério aprovado pelo Inep, e já utilizado na avaliação do desempenho linguístico nas edições do Enade desde 2013, reúne as competências acima descritas em três grupos da seguinte forma : aspectos textuais ( 40 \% ), aspectos morfossintáticos e vocabulares ( 40 \% ) e aspectos ortográficos ( 20 \% ).
A seguir, analisam -se qualitativamente os resultados da avaliação de cada um dos aspectos avaliados nos textos de resposta das questões discursivas 1 e 2 de Formação Geral.
Aspectos textuais Esta competência é a que se revelou como a mais problemática entre os participantes, porque são muitos os desvios observados.
Em relação à estrutura formal do texto produzido pelos participantes, observaram -se relações linguísticas quase agramaticais, como as estabelecidas pela sequência de gerúndios sem o apoio de um ponto de partida para a organização das informações gramaticais e semânticas, ou seja, sem oração principal.
Observou -se que uma grande parte dos participantes não organizou as ideias em parágrafos, talvez devido ao pequeno número de linhas disponibilizadas para a resposta da questão ou, quem sabe, pela suposição de que não seria necessária essa divisão por não se tratar de um texto no modelo de uma redação dissertativo-argumentativa, como solicitado nos vestibulares.
Quanto à utilização dos mecanismos de referenciação, deve -se destacar a ocorrência, em uma parte dos textos, de repetições de palavras ou expressões sem a utilização de termos sinônimos ou pronomes, como seria adequado.
Outro aspecto relevante a ser destacado é a quase total ausência de operadores argumentativos, tanto intrafrasais como interfrasais, repercutindo uma tendência atual da mídia escrita.
A respeito do uso dos sinais de pontuação, observaram -se vários desvios nos textos analisados, com a ocorrência, inclusive, de parágrafos sem marca interna de pontuação para separar os períodos.
São os seguintes os tipos de problemas encontrados, que se caracterizam como desvios ao estabelecido pela norma-padrão : a ) uso inadequado de vírgula : • ocorrência entre o sujeito e o predicado, o que é considerado inadequado por interromper a sequência natural da frase ; • ocorrência de apenas uma das vírgulas para intercalar uma palavra, uma expressão ou uma oração encaixada, quando o correto seria a demarcação com uma vírgula antes e uma depois do termo encaixado na frase ; • ocorrência da vírgula no lugar do ponto para separar ideias que constituem períodos distintos ; b ) uso inadequado de ponto e vírgula para separar elementos que deveriam ser separados por vírgula, por estarem dentro do mesmo período ; c)tausência de ponto final para separar períodos.
Vale observar que não foi penalizada a ausência de vírgula para destacar locuções ou adjuntos adverbiais de pequena extensão deslocados de posição na frase, por ser um uso opcional.
Os problemas relativos à estruturação textual geram, portanto, os seguintes tipos de desvios : • sequência justaposta de ideias sem encaixamentos sintáticos ; • redução de estruturas subordinadas, ao lado do aumento na frequência de estruturas coordenadas e absolutas ; • redução no uso de conectores para expressar relações lógicas essenciais à construção do texto ; • emprego inadequado de operadores que não estabelecem relações lógicas coerentes entre ideias do texto ; • repetição de termos sem a utilização de procedimentos mais sofisticados de substituição ( hiperonímias, hiponímias, nominalizações, expressões metafóricas ) ; • frases fragmentadas que comprometem a estrutura lógico-gramatical ; • frases formadas apenas por oração subordinada, sem oração principal.
Aspectos morfossintáticos e vocabulares Os resultados relativos aos aspectos morfossintáticos ( concordância, regência, colocação de pronomes, flexão nominal e verbal ) evidenciam desvios recorrentes no desempenho dos participantes, relacionados ao seu processo de formação e, principalmente, a hábitos da oralidade, marcados pela informalidade.
Apesar de exemplificarem processos de mudança em curso na Língua Portuguesa, esses desvios foram penalizados por não atenderem às exigências da norma-padrão relativas ao texto escrito em situação formal, como deveria ser o texto dissertativo solicitado no Enade.
O desvio mais frequente, em relação à regência, é a falta do sinal indicativo da crase – isso revela que o usuário não tem consciência de que, sob a forma do termo “ a ”, existe a presença de uma preposição “ a ”, exigida pela regência do termo anterior, combinada a um artigo definido ( por exemplo : “ Devido as grandes guerras que estão acontecendo... ” ) Outros problemas relacionados à regência verbal e à nominal foram identificados em alguns textos dos participantes : ( a ) a ausência de preposição ou seu uso indevido antes de pronome relativo ( por exemplo : “ a região que eles moram não oferece condições de segurança ” ) ; ( b ) o uso indevido do pronome relativo ( por exemplo : “ não pensam na cidade para que eles vão ” ) ; ( c ) utilização de pronome “ onde ” deslocado de sua função locativa ).
Esses processos refletem hábitos da modalidade oral da língua, em situações de registro informal.
Apesar da possibilidade de que essa alteração de regência se generalize no padrão escrito da Língua Portuguesa, como já está ocorrendo até em textos jornalísticos, a ausência da preposição foi penalizada neste processo de avaliação.
Outro desvio observado nos textos diz respeito aos processos de concordância verbal e de concordância nominal.
Quanto à concordância de número, além das ocorrências já conhecidas de falta da marca de plural em verbos ou adjetivos, em função da ocorrência de substantivo plural na função de sujeito, observaram -se dois tipos de desvio já registrados nas pesquisas linguísticas : ( a ) a ausência de marca de plural com sujeito posposto ( por exemplo : “ Foi expulso do país os imigrantes que entraram sem autorização da polícia ” ) ; ( b ) o uso indevido da marca de plural para estabelecer concordância com o plural da locução adjetiva, apesar de o substantivo que funciona como núcleo do sintagma nominal estar no singular ( por exemplo : “ A proteção das mulheres que sofrem violência foram garantidas pela Lei Maria da Penha ” ).
Uma ocorrência generalizada foi a ausência de acento circunflexo na forma plural do presente do indicativo do verbo “ ter ”, que foi considerada como um desvio na concordância verbal e não meramente de acentuação gráfica.
Quanto à concordância de gênero, alguns casos foram observados, normalmente no âmbito de sintagmas nominais longos, em que o adjetivo está afastado do substantivo.
Quanto à questão da colocação pronominal, foram poucos os desvios observados.
Apesar de serem aspectos relacionados à oralidade, observou -se que a maioria dos participantes já incorporou regras como a não introdução da frase por um pronome oblíquo e a próclise na presença de um termo atrator.
Não se adotou, entretanto, o padrão excessivamente formal descrito pelas gramáticas normativas em relação à posição do pronome oblíquo em locuções verbais, já que esse uso está muito distante da prática cotidiana, até em textos mais formais. relação à posição do pronome oblíquo em locuções verbais, já que esse uso está muito distante da prática cotidiana, até em textos mais formais.
Quanto aos aspectos vocabulares ( precisão semântica do vocabulário utilizado e adequação ao registro formal exigido pelo tipo de texto solicitado no enunciado ), podem ser destacados os seguintes tipos de desvios : expressões da oralidade, seleção vocabular incompatível com o contexto, gerando situações de falta de inteligibilidade ; falta de domínio de vocabulário mais abstrato e de maior complexidade, essencial ao desenvolvimento do texto de base dissertativa.
Não se observou presença de palavras de baixa calão, ao contrário do exame de 2015, em que o tema “ funk ” gerou a utilização desse tipo de vocabulário.
Várias marcas de oralidade foram identificadas, embora não em alta frequência : o uso do pronome relativo “ onde ” como relativo universal, falta de artigo definido antes de substantivo, repetição de palavras por falta de vocabulário, reduções como “ tá ”, “ pra ”, “ pro ”, “ prum ”, expressões informais.
Aspectos ortográficos : O desempenho dos participantes revelou uma diferença muito grande nos dois aspectos de caráter ortográfico analisados neste item : grafia e acentuação.
Observou -se a ocorrência de baixo índice de desvios na grafia das palavras e grande índice de desvios no uso dos sinais de acentuação gráfica.
Em vários casos, ocorre ausência completa de acentuação.
Os resultados revelam que existe uma tendência crescente entre os universitários brasileiros de eliminação da acentuação gráfica, talvez motivada pelos hábitos relacionados às redes sociais ou devido a um conhecimento ainda superficial sobre o Acordo Ortográfico da Língua Portuguesa de 1990.
Os casos mais sistemáticos de eliminação do acento indicador da sílaba tônica são : • palavras proparoxítonas ( por exemplo : “ proximos ”, “ politicos ”, “ publica ”, “ fisica ”, “ psicologica ” ) ; • palavras paroxítonas terminadas em ditongo crescente ( por exemplo : “ necessario ”, “ noticiarios ”, “ individuo ”, “ dependencia ”, “ varios ”, “ propria ”, “ violencia ”, “ ocorrencia ”, “ suicidio ” ) ; • palavras oxítonas ( por exemplo : “ ate ”, “ tambem ”, “ mes ” ) ; • palavras paroxítonas terminadas em ditongo nasal ( por exemplo : “ orgão ”, “ orfão ” ).
Quanto ao domínio das convenções relativas à grafia das palavras, pode -se concluir que os desvios estão majoritariamente relacionados a questões relativas à oralidade : • hipercorreção pela escolha de “ e ” no lugar de “ i ” ou de “ i ” no lugar de “ e ”, por influência de hábitos da oralidade : por exemplo, “ inviado ” ( no lugar de “ enviado ” ), “ entervenção ” ( no lugar de “ intervenção ” ), “ empedimento ” ( no lugar de “ impedimento ) ; • eliminação do “ r ” marcador do infinitivo verbal : por exemplo, “ esta ” ( no lugar de “ estar ” ), “ estuda ” ( no lugar de “ estudar ” ), “ migra ” ( no lugar de “ migrar ” ) ; • utilização de “ s ” no lugar de “ c ” : por exemplo, “ mensionado ” ( no lugar de “ mencionado ” ) ; “ insentivou ” ( no lugar de “ incentivou ” ) ; • utilização de “ s ” no lugar de “ x ” : por exemplo, “ espectativa ” ( no lugar de “ expectativa ”, devido à confusão com o verbo “ esperar ” ) ; e a situação contrária, “ x ” no lugar de “ s ”, “ expantoso ” ( no lugar de “ espantoso ” ) ; • utilização de “ ss ” no lugar de “ ç ” : por exemplo, “ gerassões ” ( no lugar de “ gerações ” ) ; • uso de “ am ” no lugar de “ ão ” na forma de terceira pessoa do singular do presente e do pretérito perfeito do indicativo : por exemplo, “ estam ” ( no lugar de “ estão ” ) ; • utilização de “ l ” no lugar de “ u ” em final de sílaba : por exemplo, “ tralma ” ( no lugar de “ trauma ” ) ; • outros casos de desvios de grafia relacionados à variação diastrática podem ser observados em “ estrupo ” ( no lugar de “ estupro ” ), “ subjulgadas ” ( no lugar de “ subjugadas ” ), “ precoseito ” ( no lugar de “ preconceito ” ), “ soubre ” ( no lugar de “ soube ” ), “ vecendo ” ( no lugar de “ vencendo ” ), “ indesencia ” ( no lugar de “ indecência ” ), “ apolojia ” ( no lugar de “ apologia ” ), “ fulga ” ( no lugar de “ fuga ” ). • inadequação no uso da maiúscula para destacar determinadas palavras-chave do texto, como “ Violência ”, “ Brasileiros ”.
Destaque -se, também, a ocorrência de participantes que grafam os textos inteiramente em caixa alta.
Vale observar que, ao contrário do esperado, não apareceram abreviaturas próprias do “ internetês ”, relacionadas ao uso de redes sociais e de e-mails.
Sua influência, portanto, parece se restringir à ausência dos acentos gráficos.
Em relação aos aspectos textuais, observa -se dificuldade na estruturação textual, o que resulta em truncamentos sintáticos e prejuízo da compreensão do texto, e desrespeito a regras essenciais de pontuação, principalmente pela dificuldade em identificar as relações lógicas que se estabelecem entre as unidades frasais.
Em relação aos aspectos morfossintáticos, domínio das regras de concordância nominal e verbal, mas desrespeito às regras de regência nominal e verbal, principalmente no que diz respeito ao uso do sinal indicativo da crase.
Em relação aos aspectos vocabulares, reduzida utilização de palavras de uso mais formal e abstrato, adequado à construção do texto de base dissertativa.
Em relação às convenções ortográficas, domínio da grafia das palavras, mas falta de atendimento à maioria das regras de acentuação gráfica. 3.3.2 - Componente de Conhecimento Específico Na parte da prova relativa às questões discursivas do Componente de Conhecimento Específico ( Tabela 3.17 ), observa -se que a Média foi bem mais baixa do que para as questões discursivas do Componente de Formação Geral.
Enquanto no Componente de Formação Geral a Média para estudantes de Tecnologia em Estética e Cosmética de todo o Brasil foi 48,5, na parte de Conhecimento Específico a Média foi 27,6.
A maior Média deste componente foi obtida pelos estudantes da região Sul ( 29,5 ), e a menor, pelos da região Norte ( 21,4 ).
Quanto à variabilidade das notas, o Desvio padrão do Brasil foi 16,3.
O maior Desvio padrão foi encontrado na região Nordeste ( 16,6 ), e o menor, na região Norte ( 15,1 ).
A maior nota Máxima, 93,3, foi obtida na região Sul.
A nota Mínima ( 0,0 ) foi obtida por pelo menos um aluno em todas as regiões do Brasil.
A Mediana do Brasil como um todo foi 26,7.
A maior Mediana foi obtida na região Sul ( 28,3 ), e a menor, na região Norte ( 20,0 ).
O Gráfico 3.10 representa a distribuição das notas nas questões discursivas no Componente de Conhecimento Específico.
A distribuição é unimodal, com moda ocorrendo no intervalo ( 20 ; 30 ].
A análise de cada uma destas questões será feita a seguir. 3.3.2.1 - Análise da Questão Discursiva 3 do Componente de Conhecimento Específico A questão 3, cujos resultados aferidos encontram -se descritos na Tabela 3.18, foi a de pior desempenho dentre as três discursivas de Conhecimentos Específicos.
A Média dos estudantes de todo o Brasil foi 19,0.
A menor Média nessa questão foi obtida pelos alunos da região Norte ( 12,8 ), enquanto a maior Média foi obtida na região Sul ( 20,6 ).
Quanto à variabilidade das notas, o Desvio padrão do Brasil foi 17,9.
O maior Desvio padrão foi obtido na região Nordeste ( 19,1 ), enquanto o menor foi obtido na região Norte ( 14,9 ).
A nota Máxima, 95,0 pontos, foi alcançada por, pelo menos, um aluno da região Sul.
A Mediana do Brasil como um todo foi 15,0, a mesma das regiões Sudeste e Centro-Oeste.
A nota Mínima ( 0,0 ) foi obtida por, pelo menos, um estudante em todas as regiões do Brasil.
O Gráfico 3.11 mostra a distribuição das notas na questão discursiva 3, do Componente de Conhecimento Específico da área de Tecnologia em Estética e Cosmética.
A moda dessa distribuição ocorre no grupo daqueles que obtiveram nota zero, dentre aqueles que tentaram responder à questão.
Essa distribuição revela uma concentração bem alta nas notas mais baixas. 3.3.2.2 - Comentários sobre as respostas à Questão Discursiva 3 A questão abordava a fisiopatologia e tratamento para as estrias, sendo um tema extremante pertinente e incluso nas competências mínimas para a formação do tecnólogo em Estética, tendo em vista que tal disfunção é bastante encontrada em nossa clínica.
O enunciado inicia -se com um texto que embasará as perguntas subsequentes de forma clara e objetiva.
Entretanto, ocorreu um grau muito acentuado de respostas fracas devido à dificuldade de interpretação da questão, no aspecto do que são teorias e fisiopatologia.
O tempo de prova e o espaço para as respostas foram suficientes, sendo a questão com um grau de dificuldade mediano.
Embora a resposta do item ‘ a ’ fosse objetiva e sem possibilidade de variação, houve uma grande quantidade de respostas erradas.
Muitos concluintes utilizaram causas de estrias como fundamento da resposta e não a teoria, o que não se pode considerar como correto.
Para o item ‘ b ’, apesar de demonstrarem ter compreendido o que se pedia, as respostas foram redigidas em linguagem bastante coloquial.
No item ‘ c ’, apesar de a pergunta ser objetiva, verificou -se bastante dificuldade para exemplificar o mecanismo de ação da microgalvanupuntura.
A pluralidade de respostas incorretas, totalizando aproximadamente 73 \%, demonstrou o fraco desempenho dos alunos de uma forma geral, não obstante tal questão ter abordado conteúdo que envolve competência básica do profissional de estética.
Cerca de 23 \% das dissertações podem ser consideradas médias e apenas 4 \%, boas.
Embora o enunciado fosse bastante claro, percebeu -se uma incapacidade de interpretação e de entendimento do que se esperava como resposta.
As respostas dos estudantes para o item ‘ a ’ revelaram que não houve entendimento do comando da questão.
Pedia -se a fisiopatologia da estria, o que foi entendido como apresentar a causa.
O mesmo tipo de erro de compreensão do comando se repetiu nas respostas dadas ao item ‘ c ’.
A dificuldade de interpretação demonstrada pelos respondentes ficou clara, nessa e nas demais questões, o que sugere uma falha na formação geral da maior parte do grupo de estudantes avaliados nessa primeira edição do Enade para a Área de Estética e Cosmética.
Em relação à parte técnica, também se notou uma disparidade ao retratar a fisiopatologia da disfunção e do mecanismo de ação do tratamento proposto.
Como tais conhecimentos são de extrema importância para a prática clínica, os cursos precisarão fazer uma análise das competências acadêmicas básicas.
Destaca -se, ainda, como importante, trabalhar com o alunado as formas de responder a perguntas como as apresentadas na parte discursiva da prova.
Foram várias as respostas sem separação das partes correspondentes aos subitens.
Algumas, apresentadas em um texto corrido, sem sequer demonstrar ordenação do pensamento para atender ao que se pedia nos subitens ‘ a ’, ‘ b ’ e ‘ c ’, o que dificultou bastante a correção dos textos. 3.3.2.3 - Análise da Questão Discursiva 4 do Componente de Conhecimento Específico A Tabela 3.19 contém as informações relativas à questão 4 do conjunto de questões do Componente de Conhecimento Específico.
O desempenho dos estudantes nessa questão foi superior ao das questões 3 e 5.
A Média geral do Brasil foi 35,9, sendo a menor Média registrada na região Norte ( 26,8 ), e a maior, na região Sul ( 37,9 ).
A nota Máxima ( 100,0 ) foi atingida por, pelo menos, um concluinte de todas as regiões.
A Mediana em todo o Brasil foi 50,0, o mesmo valor foi obtido nas regiões Sudeste, Sul e Centro-Oeste.
Nas regiões Norte e Nordeste, a Mediana foi menor ( 25,0 ).
As notas mínimas foram zero em todas as regiões.
O Gráfico 3.12 representa a distribuição de notas da questão discursiva 4, do Componente de Conhecimento Específico.
A maior frequência dessa distribuição ocorre no intervalo ( 40 ; 50 ].
O intervalo ( 20 ; 30 ] e o grupo daqueles que obtiveram zero se configuram como máximos locais destacados. 3.3.2.4 - Comentários sobre as respostas à Questão Discursiva 4 A questão 4 abordou a disfunção da acne e seu tratamento.
Essa questão, a princípio, era a mais fácil das três questões discursivas, pois abordava uma competência essencial para a formação profissional esteticista.
O enunciado iniciava -se com um estudo de caso que apresentava tanto o caso quanto o tratamento proposto.
Com o intuito de ilustrar a questão e, possivelmente, facilitar a elaboração da resposta, foi apresentada uma figura com a fisiopatologia da disfunção.
O enunciado era claro e objetivo, sendo que o grau de dificuldade da questão era baixo.
Contudo, percebe -se a dificuldade no entendimento da expressão fisiopatologia que foi respondido com graus da acne.
O tempo de prova e os espaços para as respostas foram suficientes.
O comando do item ‘ a ’ era objetivo e claro ao tratar da fisiopatologia da acne.
Contudo, de forma geral, o alunado apresentou como resposta os graus de acne, gerando um grande percentual de erros.
Como na questão 3, ficou evidenciada a dificuldade de interpretação do comando da questão.
Os estudantes não apresentaram a fisiopatologia, como foi solicitado, mas os graus da acne, mesmo havendo uma figura ilustrativa que poderia \textbf{ajudar} a responder adequadamente.
O comando do item ‘ b ’ solicitava que a resposta contivesse duas partes : a indicação da limpeza de pele e o uso do ácido mandélico.
Como não foi mencionada a quantidade de justificativas que o respondente deveria apresentar para uso desse tratamento, a indicação de uma justificativa foi considerada como correta.
Também foi considerado como acerto da questão uma mesma resposta podendo atender aos dois subitens, o que ocorreu em algumas provas.
Um aspecto que merece destaque foi a não concordância de muitos estudantes em relação ao uso do ácido mandélico, juntamente com a indicação do salicílico, o que não respondia ao que foi solicitado.
Nesse item, como se exigiu apenas uma indicação para cada item, o grau de acerto foi mais elevado.
Embora esta tenha sido a questão que apresentou melhor desempenho, por ter grau de dificuldade baixo, de uma forma geral, podemos classificar a performance como mediana.
O fato de aceitar como resposta correta e completa a apresentação de apenas uma indicação para o uso da limpeza de pele e do ácido mandélico, resultou no aumento da média da questão.
Da mesma forma que descrita na análise dos resultados da questão 3, ficou clara a necessidade de trabalhar leitura e interpretação com os estudantes, pois se percebeu que o não entendimento do comando foi a principal causa de erros. 3.3.2.5 - Análise da Questão Discursiva 5 do Componente de Conhecimento Específico A Tabela 3.20 contém as informações relativas à questão 5 do conjunto do Componente de Conhecimento Específico.
O desempenho dos estudantes de todo o Brasil nesta questão foi mais alto do que na questão discursivas 3 desse componente e mais baixo do que na questão discursiva 4.
A nota Média dos estudantes de todo o Brasil foi 27,9.
A maior Média foi registrada na região Nordeste ( 31,2 ), enquanto a menor Média foi registrada na região Centro-Oeste ( 24,4 ).
Quanto à variabilidade das notas, o Desvio padrão dos alunos do Brasil, como um todo, foi 24,5.
Enquanto o maior desvio foi encontrado na região Nordeste ( 25,1 ), o menor foi encontrado na região Norte ( 22,2 ).
A Mediana do Brasil foi 20,0 e o mesmo valor ocorreu também nas regiões Norte, Sudeste e Centro-Oeste.
Nas regiões Nordeste e Sul, a Mediana foi 25,0.
A nota Mínima ( 0,0 ) foi a mesma em todas as regiões.
Para o conjunto de alunos de Tecnologia em Estética e Cosmética do Brasil, a nota Máxima foi 100,0, tendo sido obtida por, pelo menos, um estudante da região Nordeste.
O Gráfico 3.13 representa a distribuição das notas da questão discursiva 5 do Componente de Conhecimento Específico.
Essa distribuição tem moda no intervalo ( 10 ; 20 ].
Destaca -se também, o grande número de estudantes que deixaram esta questão em branco e o intervalo ( 60 ; 70 ], com cerca de 13 \% do total de respostas, caracterizando máximos locais. 3.3.2.6 - Comentários sobre as respostas à Questão Discursiva 5 O tema abordado na questão – alopecia – faz parte das competências do profissional de estética.
Contudo, a alopecia de tração não é a abordagem mais comum do tema, fazendo com que o grau de dificuldade da questão seja considerado mais alto.
Da mesma forma que ocorreu nas questões discursivas 3 e 4, ainda que o enunciado fosse claro e objetivo, os respondentes demonstraram dificuldade no entendimento do comando.
Em relação ao item ‘ a ’, foi comum identificarem as fases, porém sem as descreverem, como solicitado : “ Identifique e descreva ”.
A maioria das respostas ao item ‘ b ’ não discorria sobre a fisiopatologia.
Para o item ‘ c ’, em raras respostas, havia justificativa para o uso de princípios ativos vasodilatadores, conforme era pedido no comando.
Em comparação com as demais questões discursivas, essa apresentava um grau de dificuldade maior.
No entanto, isso não se traduziu nas notas obtidas, pois o desempenho ficou no mesmo patamar das demais discursivas.
O tempo e os espaços foram condizentes.
O item ‘ a ’ solicitava a identificação e a descrição do ciclo do cabelo, mas, em sua maioria, as respostas traziam apenas identificação, sem qualquer descrição do ciclo.
Foi comum identificarem as fases do ciclo do cabelo, porém sem descrevê -las.
O comando do item ‘ b ’ era claro, contudo percebeu -se dificuldade de compreensão do termo fisiopatologia, o que levou a um percentual elevado dos erros nesse item.
As respostas dadas ao item ‘ c ’ demonstraram o não conhecimento da justificativa e da associação do microagulhamento.
O desempenho dos alunos, de forma geral, foi fraco.
A média foi considerada baixa, apesar de essa ser a questão mais complexa dentre as três discursivas.
A temática da terapia capilar parece ser mais desconhecida, já que essa também foi a questão com maior índice de respostas em branco.
Com respeito à formação técnica, percebe -se dificuldade em relação à fisiopatologia da disfunção tratada nessa questão, como nas duas outras.
Os currículos deveriam dar uma atenção especial ao segmento de terapia capilar que parece não ter sido explorado de forma mais significativa na formação dos concluintes participantes do Enade/2016, embora seja uma competência presente nos currículos do curso. 3.3.3 - Considerações Finais É importante ressaltar que os cursos de Tecnologia em Estética e Cosmética entraram no catálogo dos cursos superiores em tecnologia em 2015 e, ainda, não há Diretrizes Curriculares específicas para eles.
Contudo, a prova foi condizente com a formação do profissional tecnólogo em Estética e Cosmética, com questões bem formuladas e pertinentes à qualificação que se espera de um profissional dessa área.

% matched lemmas: ajudar
\end{textsample}
